% uid: ""
% school: "lycee"
% level: "3eme"
% section: ["Economie et Gestion"]
% chapter: "Statistiques"
% title: "Prix et demande : ajustement et prévision"
% difficulty: "Moyen"
% points: null
% tags: ["statistiques", "serie-double", "droite-de-mayer", "demande", "prevision", "interpretation"]
% provenance: "own"
% dressing: "ajustement affine"
% rung: 3
% batch: "2026-11-20-stat-3eco-A"

% ---STATEMENT---
Une coopérative a observé la quantité demandée $y$ (en tonnes) d'un produit selon son prix $x$ (en dinars le kilo) :

\begin{center}
\begin{tabular}{|l|c|c|c|c|c|c|c|c|}
\hline
Prix $x_i$ & $2$ & $3$ & $4$ & $5$ & $6$ & $7$ & $8$ & $9$ \\ \hline
Demande $y_i$ & $96$ & $90$ & $83$ & $78$ & $70$ & $64$ & $57$ & $52$ \\ \hline
\end{tabular}
\end{center}

\begin{enumerate}
  \item Calculer les points moyens $G_1$ (quatre premiers prix) et $G_2$ (quatre derniers).
  \item Déterminer l'équation de la droite de Mayer $y=ax+b$. Interpréter le signe de $a$.
  \item Estimer la demande si le prix est fixé à $10$ dinars.
  \item Pour quel prix cet ajustement prévoit-il une demande nulle ? Ce résultat est-il utilisable pour un prix de $20$ dinars ? Commenter.
\end{enumerate}

% ---SOLUTION---
\begin{enumerate}
  \item Points moyens des deux sous-nuages de $4$ points (prix $2$ à $5$, puis $6$ à $9$) :
  $$G_1:\ \overline{x}_1=\frac{2+3+4+5}4=3{,}5,\quad\overline{y}_1=\frac{96+90+83+78}4=\frac{347}4=86{,}75.$$
  $$G_2:\ \overline{x}_2=\frac{6+7+8+9}4=7{,}5,\quad\overline{y}_2=\frac{70+64+57+52}4=\frac{243}4=60{,}75.$$
  \item Pente : $a=\dfrac{60{,}75-86{,}75}{7{,}5-3{,}5}=\dfrac{-26}4=-6{,}5$. Ordonnée à l'origine, avec $G_1$ : $86{,}75=-6{,}5\times3{,}5+b\Rightarrow b=86{,}75+22{,}75=109{,}5$.
  $$ \ib{$y=-6{,}5\,x+109{,}5$}.$$
  $a<0$ : quand le prix augmente de $1$ dinar, la demande baisse d'environ $6{,}5$ tonnes (loi de la demande).
  \item $x=10$ : $y=-6{,}5\times10+109{,}5=-65+109{,}5=44{,}5$ tonnes.
  \item $-6{,}5\,x+109{,}5=0\iff x=\dfrac{109{,}5}{6{,}5}=\dfrac{219}{13}\approx16{,}85$ dinars.

  Pour $x=20$ : $y=-130+109{,}5=-20{,}5<0$ : une quantité ne peut pas être négative.

  \textbf{Piège :} l'ajustement n'est valable que sur (ou près de) l'intervalle observé $2\le x\le9$. Au-delà de $16{,}85$ dinars, la droite perd tout sens économique : il ne faut jamais extrapoler loin des données.
\end{enumerate}

\textbf{Idée :} un ajustement affine décrit la tendance dans la zone observée, pas au-delà.
